\chapter{Introduction}

\section{Background}

Tensors higher dimensional generalisations of vectors and matrices.
Tensor neworks and tensor contraction, which will be explained in detail in chapter \ref{ch:tensors}, are ubiquitous in computer science, physics, and applied mathematics.
This thesis focuses on the practical problem of finding an efficient tensor contraction algorithm.

Most existing tensor contraction algorithms are focused on find an optimal contraction order of tensors' bonds.
As we will show in the chapter \ref{ch:treewidth}, finding the optimal order is NP-complete and is equivalent to finding the treewidth of the tensor network's line graph.

\section{Outline}

We start this thesis by a thorough introduction to tensors and tensor networks in chapter \ref{ch:tensors}.
In chapter \ref{ch:treewidth}, we will introduce the concept of treewidth and how treewidth can bu used to determine the complexity of tensor contraction.
Chapter \ref{ch:zx_calculus} describes a popular graphical language, ZX calculus, for a special type of tensor networks called quantum circuits. 
The next chapter \ref{ch:zx-algorithms} combines the idea of ZX calulus and treewidth to propose a novel tensor contraction algorithm for quantum circuits.
In the last chapter \ref{ch:generic_diagram}, we will discuss several other diagrammatic calculus that w

\section{Statement of Originality}

Most of the results in this thesis are well established in literature and they are cited properly. 
All of the explanations, proofs, diagrams, demonstrations, and discussions are my own work and interpretations, unless otherwise stated.

All of the proofs in chapter \ref{ch:treewidth} are my own, although they are classical results in graph theory.
They drew many inspirations from \textit{Introduction to Graph Theory} by Douglas B. West \cite{west2001introduction} and several other papers including
\cite{dirac1961rigid,fulkerson1965incidence,rose1970triangulated,gavril1974intersection}.

The section on ZX and ZXW calculus in chapter \ref{ch:zx_calculus} and \ref{ch:zx-algorithms} by mostly based on the book by Kissinger and van de Wetering \cite{KissingerWetering2024Book}, as well as several related papers including those by Wang \cite{wang2025penrosetensordiagramszx, Wang_complete}.

The tensor contraction algorithms using treewidth and local complementations are novel, although it was inspired by the paper by Cam and Martiel \cite{cam2023speedingquantumcircuitssimulation}.
The ideas of generalised complete tensor notation is original, although it definitely has been discussed in the literature.

Another important part of this thesis is the program I wrote to implement the algorithms. 
This program was written in Rust and also includes a generic graph theory library, a tensor contraction library, and an online interactive graph board.

The figures in this thesis drew inspiration from the illustrations in \cite{KissingerWetering2024Book} and \cite{waintal2026competequantumcomputerslecture}. Most of them are, however, drawn by myself using tools I created.

I use open-sourced LLM by DeepSeek and proprietory LLM ChatGPT to help with proofreading, literature search, and some coding tasks.
No other parts of this thesis were involved with any AI tools. 
No other AI models are used.

The thesis is typesetted using XeLateX using CUED Template created by \href{https://github.com/kks32/phd-thesis-template} {Krishna Kumar}. 
The figures are created using TikZ and my own graphboard program at \url{https://github.com/harryhanYuhao/graphboard}.
